\section{Appendix D: Zeckendorf Representation and Stabilization Kernel Properties}
\label{app:zeckendorf-fold}

\subsection{Counting Recursion Proof}

Let $a_m=\card{X_m^{\mathrm Z}}$. Classify by final bit:
\begin{itemize}
  \item Final bit $0$: first $m-1$ bits arbitrary legal, count is $a_{m-1}$;
  \item Final bit $1$: second-to-last bit must be $0$, first $m-2$ bits arbitrary legal, count is $a_{m-2}$.
\end{itemize}
Thus $a_m=a_{m-1}+a_{m-2}$, which combined with $a_1=2,a_2=3$ gives $a_m=F_{m+2}$.

\subsection{\texorpdfstring{Idempotence and Surjectivity of $\mathrm{Fold}_m$}{Idempotence and Surjectivity of Fold\_m}}

\textbf{Idempotence:} If $x\in X_m^{\mathrm Z}$, then its Fibonacci weighted value is already a Zeckendorf legal representation, so folding does not change it, i.e., $\mathrm{Fold}_m(x)=x$.\\
\textbf{Surjectivity:} For any legal $x$, take the integer $N=\sum_{k=1}^m x_k F_{k+1}$, whose Zeckendorf first $m$ bits are $x$, hence there exists $\omega$ such that $\mathrm{Fold}_m(\omega)=x$.
